% Integrated from ns_arithmetic_flow.tex SHA256
% 82e369b60c195d06d3bc4a07b47e2b2ec764c1b9e36bc3a5982179ff87cca7cc.
% Full required Fourier-domain proofs from heat_connection.tex and
% xi4_arithmetic_transport.tex are retained below; exact source hashes
% and label maps are recorded in the integration receipt.
\section{Exact arithmetic profiles driven by a supplied incompressible flow}
\label{sec:ns-full-flow}

\subsection{Public source, retained input class and coordinates}

The public manuscript \emph{Finite Time Blowup for Navier--Stokes}
\cite[Theorem 1.1]{openaiNS2026release} claims a smooth forced velocity on
$\R^3\times[0,1)$, with one fixed compact spatial support, initial
velocity zero, bounded kinetic energy and unbounded terminal supremum
norm, for every $\nu>0$. Its complete proof and formal verification are
not imported as independently certified results in this section. We
construct and prove an operator correspondence for its stated smooth
input class. Section~\ref{sec:ns-actual-witness} applies the complete
correspondence to the source's actual cutoff-summed corrected field.
The correspondence itself is established
below without assuming any unproved endpoint. The first source statement
is an attribution; the formulas and estimates that follow are our proofs.

For attribution, Buckmaster's primary statement
\cite[p.~1]{buckmasterStatement2026} identifies his collaboration with
Levent Alp\"oge on smooth-forced Euler and credits Diego C\'ordoba and
Luis Mart\'inez-Zoroa for the underlying forced-blowup programme.
The revised NS manuscript \cite[pp.~2--3]{openaiNS2026release} now discusses
that programme and the hypodissipative extension with Zheng explicitly.
These statements identify the respective work. No dependence of that
historical programme on the polynomial coordinates below is asserted.

Let $\mathcal U_{K,T,\nu}$ be the set of smooth real triples $(u,p,f)$
on $[0,T)\times\R^3$ satisfying, component by component,
\begin{equation}
 \partial_tu_i+\sum_{j=1}^3u_j\partial_ju_i
 =\nu\sum_{j=1}^3\partial_j^2u_i-\partial_ip+f_i,
 \qquad \sum_i\partial_i u_i=0,
 \qquad \operatorname{supp}u(t,\cdot)\subset K,
 \label{nfl:input}
\end{equation}
where $K\subset\R^3$ is compact, $0<T<\infty$, and $\nu>0$ is retained.
We do not require compact support of $p,f$ for the following map. The
actual compactly supported source input in
Section~\ref{sec:ns-actual-witness} satisfies these conditions on every
preterminal interval. No assertion of existence of
a singular member is part of the definition of this set.

In the original real polynomial coordinates $x=(x_1,x_2,x_3)=(x,y,w)$,
retain every coefficient of $F$ and put
\begin{equation}
 S(x)=(F_1(x)/2,F_2(x),F_3(x))^{\mathsf T},\quad
 J(x)=DS(x)=\operatorname{diag}(1/2,1,1)DF(x).
 \label{nfl:S}
\end{equation}
The determinant proof in \eqref{eq:ifh-det} gives $\det J=-1$, with the exact inverse
\begin{equation}
 J^{-1}=(DF)^{-1}\operatorname{diag}(2,1,1)
       =[2W_1\ W_2\ W_3].
 \label{nfl:Jinverse}
\end{equation}
Thus $S_1=\sigma$ is precisely the earlier arithmetic coordinate. The
other two coordinates retain the entire original target. No global
injectivity of $S$ or $F$ is used; $J^{-1}$ is a pointwise matrix inverse.
Both $J$ and $J^{-1}$ have polynomial entries. Here the columns
$W_i=(DF)^{-1}e_i=-\tfrac12\nabla F_{i+1}\times\nabla F_{i+2}$,
with cyclic indices, are the full inverse columns proved in
Section~\ref{sec:inverse-fibre-heat}.
The first-coordinate observation alone has the exact velocity kernel
$\ker(d\sigma)_x=\operatorname{span}_{\R}\{W_2(x),W_3(x)\}$.
Indeed $d\sigma(W_i)=\delta_{1i}/2$ and the $W_i$ form a basis, by
the full inverse matrix. The retained two additional target coordinates
recover precisely those two directions in \eqref{nfl:Jinverse}.

\subsection{The actual arithmetic heat function}
We use the convention of \cite[Section 1, equations (1)--(4)]{rodgersTao2020}:
\begin{align}
\Phi(u)&=\sum_{n\ge1}(2\pi^2n^4e^{9u}-3\pi n^2e^{5u})
                       e^{-\pi n^2e^{4u}},\notag\\
H_t(z)&=\int_0^\infty e^{tu^2}\Phi(u)\cos(zu)\,du,\notag\\
Z_t(s)&=8H_t(2s),\qquad
Z_0(s)=\Xi(s)=\xi(1/2+is),\notag\\
\xi(\sigma)&=\tfrac12\sigma(\sigma-1)\pi^{-\sigma/2}
                     \Gamma(\sigma/2)\zeta(\sigma).
\label{nfl:k:Z}
\end{align}
The Mellin and spectral coordinates are related by
$\sigma=1/2+is$, with inverse $s=-i(\sigma-1/2)$.
They are not silently identified.

We rederive the integral identity and its analytic bounds. Set
\[
h(x)=\frac\pi2x^2(2\pi x^2-3)e^{-\pi x^2},\quad
S(\lambda)=\sum_{n\ge1}h(n\lambda),\quad
k(v)=e^{v/2}S(e^v).
\]
For the Fourier convention $\widehat h(q)=\int_\R h(x)e^{-2\pi ixq}dx$,
the Gaussian transform and its second and fourth derivatives give
\begin{align*}
\widehat{x^2e^{-\pi x^2}}&=(1/(2\pi)-q^2)e^{-\pi q^2},\\
\widehat{x^4e^{-\pi x^2}}&=(q^4-3q^2/\pi+3/(4\pi^2))e^{-\pi q^2}.
\end{align*}
Multiplying by $-3\pi/2$ and $\pi^2$ respectively gives $\widehat h=h$.
Since $h(0)=0$, Poisson summation gives
$S(\lambda)=\lambda^{-1}S(1/\lambda)$ and $k(v)=k(-v)$.
For $\lambda\ge1$,
$|S(\lambda)|\le C\lambda^4e^{-\pi\lambda^2}$, because the
remaining series is bounded by a constant times
$\sum_{n\ge1}n^4e^{-\pi(n^2-1)}<\infty$.
The symmetry gives the corresponding bound as $\lambda\downarrow0$.
Each fixed number of $v$ derivatives introduces only finitely many
additional powers of $ne^v$. It follows, with constants depending
on the derivative order, that
\begin{equation}
|k^{(j)}(v)|\le C_j e^{C_j|v|}e^{-\pi e^{2|v|}}.
\label{nfl:k:kernel-bound}
\end{equation}
Termwise differentiation on compact sets is justified by the same
summable bounds. On the negative half-line differentiate the proved
evenness identity. This proves \eqref{nfl:k:kernel-bound} on both ends.

For $\Re\sigma>1$, absolute integrability permits summation inside
the Mellin integral. The substitution $q=\pi x^2$ gives
\begin{align*}
\int_0^\infty h(x)x^{\sigma-1}dx
 &=\pi^{-\sigma/2}
   \left[\tfrac12\Gamma(\sigma/2+2)-\tfrac34\Gamma(\sigma/2+1)\right]\\
 &=\frac{\sigma(\sigma-1)}8\pi^{-\sigma/2}\Gamma(\sigma/2),\\
\int_0^\infty S(\lambda)\lambda^{\sigma-1}d\lambda&=\xi(\sigma)/4.
\end{align*}
The two-ended bounds make the last integral entire in $\sigma$,
so the equality extends by analytic continuation. Substituting
$\lambda=e^v$ and $\sigma=1/2+is$ proves
\begin{equation}
Z_t(s)=4\int_\R e^{tv^2/4}k(v)e^{isv}\,dv.
\label{nfl:k:Fourier}
\end{equation}
Indeed direct expansion also gives $\Phi(u)=2k(2u)$; the substitution
$v=2u$ and the evenness of $k$ carry the stated $H_t$ integral
to \eqref{nfl:k:Fourier}, retaining the factors $2,4,8$.
The bound \eqref{nfl:k:kernel-bound} dominates every fixed derivative
uniformly on compact complex $(t,s)$ sets. Thus $Z$ is jointly entire
and differentiation under its integral proves exactly
\begin{equation}
\partial_t Z_t=-\tfrac14\partial_s^2 Z_t.
\label{nfl:k:heat}
\end{equation}
This supplies the actual initial datum and the actual time orientation.

\subsection{An explicit domain for the heat evolution}
The differential expression alone does not specify a heat evolution
on every entire-function space. We now give a precise invariant domain
containing the actual kernel.

\begin{definition}
Let $\mathscr E_{\rm nfl}$ consist of $\phi\in C^\infty(\R,\C)$ for which
\[
p_{A,B,m,j}(\phi)=\int_\R |\phi^{(j)}(v)|(1+|v|)^m
                          e^{Av^2+B|v|}\,dv<\infty
\]
for every nonnegative integer $A,B,m,j$. Give it these seminorms.
Define
\[
T\phi(s)=4\int_\R\phi(v)e^{isv}\,dv,\qquad \mathscr A_{\rm nfl}=T\mathscr E_{\rm nfl},
\]
and equip $\mathscr A_{\rm nfl}$ with the transported topology.
\end{definition}

The map $T$ is injective. To give the needed uniqueness argument,
let $T\phi$ vanish on the real line and let
$g_\epsilon(v)=(4\pi\epsilon)^{-1/2}e^{-v^2/(4\epsilon)}$.
Its Gaussian integral representation is
\[
 g_\epsilon(v)=(2\pi)^{-1}\int_\R e^{-\epsilon s^2}e^{isv}ds.
\]
Fubini is valid because $\phi\in L^1$ and the Gaussian is integrable;
it expresses $\phi*g_\epsilon$ as a Gaussian-weighted integral of the
vanishing Fourier transform. Thus $\phi*g_\epsilon=0$.
As $\epsilon\downarrow0$, these convolutions converge to $\phi$ in
$L^1$: translations are continuous in $L^1$, first for continuous
compactly supported functions by uniform continuity, then for all
$L^1$ functions by approximation; splitting the Gaussian integral
into a small translation ball and its vanishing tail proves the claim.
Consequently $\phi=0$ almost everywhere and, being continuous, everywhere.
This also justifies the transported topology without an unspecified
choice of representative.

\begin{proposition}[Heat evolution with its inverse]
For every complex $t$, multiplication
$E_t\phi(v)=e^{tv^2/4}\phi(v)$ is a continuous automorphism of $\mathscr E_{\rm nfl}$.
The operators $U_t=TE_tT^{-1}$ on $\mathscr A_{\rm nfl}$ satisfy
\begin{gather}
U_tU_u=U_{t+u},\qquad U_t^{-1}=U_{-t},\qquad
\partial_tU_t=-\tfrac14\partial_s^2U_t,\notag\\
U_t\Xi=Z_t,\qquad U_t\partial_s=\partial_s U_t,\qquad
U_tM_sU_{-t}=M_s-\tfrac t2\partial_s. \label{nfl:k:conjugation}
\end{gather}
Here $M_s$ is multiplication by the unchanged spectral coordinate.
\end{proposition}
\begin{proof}
Every $v$ derivative of $e^{tv^2/4}\phi$ is a finite sum of
$e^{tv^2/4}$ times a polynomial in $v$ times a derivative of $\phi$.
Each target seminorm is bounded by finitely many source seminorms
with larger $A,m$; on compact $t$ sets these bounds are uniform.
The inverse multiplier is $e^{-tv^2/4}$. Differentiation in $t$
inserts $v^2/4$, while two $s$ derivatives in the transform insert
$-v^2$, proving the generator and the group identities.
The kernel bound proves $k\in\mathscr E_{\rm nfl}$, so $U_t\Xi=Z_t$ follows from
\eqref{nfl:k:Fourier}.

Integration by parts gives $M_sT\phi=T(i\phi')$.
For completeness the boundary terms vanish: each derivative of
$e^{isv}\phi$ is integrable on a horizontal real-$v$ line and the
function itself is integrable, so its limits at both ends exist and
are zero. The same argument applies after exponential weights.
Thus multiplication by $s$ preserves $\mathscr A_{\rm nfl}$; differentiation does also,
since $\partial_sT\phi=T(iv\phi)$.
Finally
\[
e^{tv^2/4}i\partial_v(e^{-tv^2/4}\phi)
   =i\phi'-\tfrac t2iv\phi.
\]
Applying $T$ proves the last identity in \eqref{nfl:k:conjugation}.
All displayed operators have the domains just specified.
\end{proof}


In particular,
\[
 k(0)=\sum_{n\ge1}\frac{\pi}{2}n^2(2\pi n^2-3)e^{-\pi n^2}>0,
\]
because each term is positive. The same exact derivative estimates imply
that $T\phi$ and every spectral derivative are Schwartz on the real line
for every $\phi\in\mathscr E_{\rm nfl}$. Indeed
$\partial_s^jT\phi=T((iv)^j\phi)$, and integration by parts $n$ times
expresses $s^n\partial_s^jT\phi$ as
$T(i^n\partial_v^n((iv)^j\phi))$. The boundary terms vanish by the
weighted integrability argument above, and the last transform is
bounded in absolute value by four times the $L^1$ norm of its
amplitude. All pairs $n,j$ are allowed, which gives every Schwartz
seminorm with no change of Fourier factor.

\subsection{An invertible arithmetic map on the full retained Fourier domain}

Retain $c_0=-59/275184$, $c_1=1/1296$, $c_2=3/132496$
from \eqref{eq:xw-source}. Define the unchanged full profile
\begin{equation}\label{nfl:m:K}
 K_\theta(s)=\mathcal K_\theta(s)
 =c_0Z_\theta(s+3)-\frac{Z_\theta'(s+3)}{336}
       +c_1Z_\theta(s+9/2)+c_2Z_\theta(s+13/2).
\end{equation}
The two symbols $K_\theta$ and $\mathcal K_\theta$ denote this
identical function throughout the present section.
Retain the exact spaces \(\mathscr E_{\rm nfl},\mathscr A_{\rm nfl}\) and transform \(T\) already
defined above, including every seminorm
\[
p_{A,B,m,j}(\phi)=\int_\R|\phi^{(j)}(v)|(1+|v|)^m
                         e^{Av^2+B|v|}\,dv,
\qquad T\phi(s)=4\int_\R\phi(v)e^{isv}\,dv.
\]
Define the source multiplier and arithmetic operator by
\begin{align}
m(v)&=c_0e^{3iv}-\frac{iv}{336}e^{3iv}
               +c_1e^{9iv/2}+c_2e^{13iv/2},
\label{nfl:m:multiplier}\\
(\mathcal T_{\mathcal D}f)(s)
 &=c_0f(s+3)-\frac{f'(s+3)}{336}
                  +c_1f(s+9/2)+c_2f(s+13/2).
\label{nfl:m:T-D}
\end{align}

\begin{lemma}[A uniform exact lower bound for the original multiplier]
\label{nfl:m:nonzero-multiplier}
For every real \(v\),
\begin{equation}
|m(v)|\geq\delta,\qquad \delta=\frac{575}{1752192}>0.
\label{nfl:m:delta}
\end{equation}
In particular its real-axis reciprocal exists everywhere, with no
branch choice or removed spectral point.
\end{lemma}

\begin{proof}
Factoring the nonzero phase gives the exact identity
\[
e^{-3iv}m(v)=c_0-\frac{iv}{336}
                  +c_1e^{3iv/2}+c_2e^{7iv/2}.
\]
For \(|v|\leq1/2\), the elementary inequality
\(\cos x\geq1-x^2/2\) yields
\begin{align*}
\operatorname{Re}(e^{-3iv}m(v))
 &=c_0+c_1\cos(3v/2)+c_2\cos(7v/2)\\
 &\geq c_0+\frac{23}{32}c_1-\frac{17}{32}c_2
 =\frac{575}{1752192}.
\end{align*}
For completeness, this cosine bound follows directly by integration:
for \(x\geq0\),
\(1-\cos x=\int_0^x\sin u\,du\leq\int_0^xu\,du\), using
\(\sin u\leq u\); the latter follows from
\(u-\sin u=\int_0^u(1-\cos w)\,dw\geq0\). Evenness gives the
negative case. Since the real part is positive, its displayed
lower bound also bounds the modulus.

For \(|v|\geq1/2\), the triangle inequality with the original
linear term gives
\begin{align*}
|m(v)|&\geq\frac{|v|}{336}-(|c_0|+c_1+c_2)\\
 &\geq\frac1{672}-\frac{10825}{10732176}
 =\frac{10291}{21464352}
 >\frac{575}{1752192}.
\end{align*}
All coefficients in these estimates are the original rational
coefficients. The two intervals cover the whole real axis and
prove \eqref{nfl:m:delta}.
\end{proof}

\begin{theorem}[Exact arithmetic automorphism and heat conjugacy]
\label{nfl:m:automorphism}
Multiplication by \(m\) is a continuous automorphism of \(\mathscr E_{\rm nfl}\).
Its inverse is literal pointwise multiplication by \(1/m\) on
the unchanged real Fourier axis. The operator
\(\mathcal T_{\mathcal D}\) is a continuous automorphism of
\(\mathscr A_{\rm nfl}\), with exact inverse
\begin{equation}
\mathcal T_{\mathcal D}^{-1}f
 =T\left(\frac{T^{-1}f}{m}\right).
\label{nfl:m:inverse}
\end{equation}
For every complex \(t\),
\begin{gather}
\mathcal T_{\mathcal D}=TM_mT^{-1},\qquad
\mathcal T_{\mathcal D}U_t=U_t\mathcal T_{\mathcal D},\qquad
\mathcal T_{\mathcal D}^{-1}U_t
        =U_t\mathcal T_{\mathcal D}^{-1},
\label{nfl:m:commutation}\\
\mathcal K_t=\mathcal T_{\mathcal D}Z_t,
\qquad
Z_t=\mathcal T_{\mathcal D}^{-1}\mathcal K_t.
\label{nfl:m:actual-inverse}
\end{gather}
Thus the actual arithmetic function is exactly recoverable from
the propagated source coefficient, on this fully specified domain.
\end{theorem}

\begin{proof}
Every derivative of \(m\) has at most linear growth on the real
axis. More explicitly, for \(j\geq0\) it satisfies
\( |m^{(j)}(v)|\leq C_j(1+|v|)\), where one may take
\[
C_j=|c_0|3^j+c_1(9/2)^j+c_2(13/2)^j
       +\frac{3^j+\mathbf1_{j\geq1}j3^{j-1}}{336}.
\]
The product rule gives the seminorm estimate
\[
p_{A,B,m,j}(m\phi)
 \leq\sum_{\ell=0}^j\binom j\ell C_\ell
                     p_{A,B,m+1,j-\ell}(\phi).
\]
This proves continuity of \(M_m\) on the original \(\mathscr E_{\rm nfl}\).

Let \(r(v)=1/m(v)\), which is smooth by
Lemma~\ref{nfl:m:nonzero-multiplier}. The exact product identity
\(mr=1\) and its derivatives give
\[
r^{(j)}=-\frac1m\sum_{\ell=1}^j
                 \binom j\ell m^{(\ell)}r^{(j-\ell)}
\quad(j\geq1),\qquad |r|\leq\delta^{-1}.
\]
Set \(Q_0=\delta^{-1}\) and recursively define the finite
positive constants
\[
Q_j=\delta^{-1}\sum_{\ell=1}^j
                      \binom j\ell C_\ell Q_{j-\ell}.
\]
Induction proves
\(|r^{(j)}(v)|\leq Q_j(1+|v|)^j\): the term with index
\(\ell\geq1\) has exponent at most
\(1+j-\ell\leq j\), and the denominator is bounded by
\(\delta^{-1}\). Therefore
\[
p_{A,B,m,j}(r\phi)
 \leq\sum_{\ell=0}^j\binom j\ell Q_\ell
                         p_{A,B,m+\ell,j-\ell}(\phi).
\]
This proves that reciprocal multiplication is continuous on
exactly the same \(\mathscr E_{\rm nfl}\). The two pointwise compositions are
the identity, proving the Fourier automorphism and its inverse.

Applying the retained transform to each term of \(m\phi\)
gives the three translations and the derivative in
\eqref{nfl:m:T-D}; the factor \(iv\) is exactly the Fourier
multiplier of \(\partial_s\). The injectivity of \(T\) already
proved above makes its inverse on \(\mathscr A_{\rm nfl}\) unambiguous. This
proves the first identity in \eqref{nfl:m:commutation} and
\eqref{nfl:m:inverse}. Multiplication by \(m\) and by \(1/m\)
each commutes pointwise with multiplication by
\(e^{tv^2/4}\). All three maps preserve \(\mathscr E_{\rm nfl}\), so applying
\(T\) proves the two full-domain heat commutation identities
for every complex \(t\). Finally \eqref{nfl:m:K} is exactly
\(\mathcal T_{\mathcal D}Z_t\), and applying the proved inverse
gives \eqref{nfl:m:actual-inverse}.
\end{proof}


\subsection{Flow, inverse and volume form on every preterminal interval}

For an input of \eqref{nfl:input}, define $X_t(q)$ by
\[
 \partial_tX_t(q)=u(t,X_t(q)),\qquad X_0(q)=q.
\]
Here $q\in\R^3$ is the retained initial label and $x=X_t(q)$ the physical
coordinate. This is a global smooth diffeomorphism for each $t<T$.
To prove existence with its actual domain, fix $T_0<T$. On
$[0,T_0]\times K$, $u$ and every spatial derivative are bounded by
compactness; they vanish outside $K$. In particular $u$ is globally
Lipschitz in space, uniformly on this time interval. On a sufficiently
short time interval, the map
$Y\mapsto q+\int_0^t u(s,Y(s))\,ds$ is a contraction in the uniform
norm on continuous curves, because its Lipschitz constant is at most
the interval length times $\sup|Du|$. Successive contractions cover
$[0,T_0]$; bounded $u$ prevents escape in each such interval. The same
argument backward in time gives the inverse flow. Differentiation of
the integral equation gives successive linear integral equations for
each label derivative, proving their existence and continuity by the
same local iteration, and proves smoothness by induction. Since $T_0$
is arbitrary this proves all the assertions on $[0,T)$.

Set $A=D_qX_t$. Differentiation gives
\[
 A'=Du(t,X_t)A,\qquad A(0)=I_3,\qquad
 (\det A)'=(\operatorname{div}u)(t,X_t)\det A=0.
\]
The determinant identity follows first near $t=0$ by the product rule
for the determinant and $A^{-1}$; it gives $\det A=1$, which prevents
loss of invertibility and continues that identity to every $t<T$.
Thus $\det A=1$ with its original positive orientation. Every point
outside $K$ has the constant trajectory, by uniqueness. If a trajectory
starting in $K$ ended outside $K$, the backward constant trajectory
from its endpoint would contradict the inverse-flow uniqueness.
Consequently
\begin{equation}
 X_t(K)=K,\qquad X_t(q)=q\ (q\notin K),\qquad
 X_t(L)\subset K\cup L\quad\hbox{for every compact }L\subset\R^3.
 \label{nfl:compact}
\end{equation}
The inverse equality follows in the same way. These conclusions require
no bound on $\sup_{t<T}|u|$.

\subsection{An invertible profile correspondence with a complex spectral coordinate}

Write $\mathcal B=C^\infty(\R^3;\mathcal O(\C_\zeta))$, with seminorms
given by uniform suprema of finitely many real $q$ derivatives and
complex $\zeta$ derivatives on compact products. Let
\[
 s_{j,t}(q)=S_j(X_t(q)),\qquad
 \mathcal P_t:\mathcal O(\C)^3\longrightarrow\mathcal B^3,
 \quad (h_j)_j\longmapsto
       \bigl(h_j(s_{j,t}(q)+\zeta)\bigr)_j.
\]
Every $s_{j,t}$ is real. The domain variables of $h_j$ are complex;
the spatial labels remain real, since a general smooth NS velocity
need not have a holomorphic extension in space. The map on the full
product is explicitly
\begin{equation}
 H_{j,t}:\R^3\times\C_\zeta\longrightarrow\R^3\times\C_s,
 \quad(q,\zeta)\mapsto(q,s_{j,t}(q)+\zeta),\qquad
 H_{j,t}^{-1}(q,s)=(q,s-s_{j,t}(q)).
 \label{nfl:skew}
\end{equation}
These are smooth in the real variables and biholomorphic on every
complex fibre. For any fixed label $q_*$, an exact left inverse of
$\mathcal P_t$, and inverse on its image, is
\begin{equation}
 (\mathcal R_tG)_j(s)=G_j(q_*,s-s_{j,t}(q_*)),\qquad
 \mathcal R_t\mathcal P_t=I.
 \label{nfl:leftinverse}
\end{equation}
Substitution proves both directions on the image. On each compact
product the image argument stays in a compact subset of $\C$, and all
the finitely many derivatives of $s_{j,t}$ involved are bounded.
The chain and product rules prove continuity of $\mathcal P_t$;
evaluation and translation prove continuity of $\mathcal R_t$.
Thus the image carries a topology for which the stated maps are an
actual topological isomorphism, namely its subspace topology in
$\mathcal B^3$. The map retains the full original complex profile.

For comparison with the earlier spatial observable, let
$e_0(q)=(q,0)$. The precise restriction is
\[
 e_0^*(\mathcal P_th)_1(q)=h_1(\sigma(X_t(q)))
     =X_t^*(\sigma^*h_1)(q).
\]
The inverse in \eqref{nfl:leftinverse} uses the complex fibre; no
holomorphic spatial flow is silently supplied to this real restriction.

\subsection{Full first and second time generators, and Euclidean diffusion}

Let $\vartheta:[0,T)\to\R$ be smooth. It is the chosen Newman-time
map, independent of physical $t$; no value for it is inferred from NS.
Take either $h_\theta=Z_\theta$ or the full four-term coefficient
$h_\theta=K_\theta$ of Section~\ref{sec:xi4-actual-Weil}, so that
$\partial_\theta h_\theta=-\partial_s^2h_\theta/4$ with its original
factor and sign. Set
\[
 \Psi_j(t,q,\zeta)=h_{\vartheta(t)}(s_{j,t}(q)+\zeta),\qquad
 v_j(t,q)=[\nabla S_j\cdot u](t,X_t(q)).
\]
Then $\partial_ts_{j,t}=v_j$ and the chain rule proves the exact equation
\begin{equation}
 \partial_t\Psi_j=-\frac{\vartheta'}4\partial_\zeta^2\Psi_j
                      +v_j\partial_\zeta\Psi_j.
 \label{nfl:first}
\end{equation}
The exact connection on the moving image of $\mathcal P_t$ is
$\nabla_t=\partial_t-\operatorname{diag}(v_j)\partial_\zeta$:
$\nabla_t(\mathcal P_t h_t)=\mathcal P_t(\partial_t h_t)$.
This follows by the same chain rule for every smooth curve of entire
profiles. It is a connection between the displayed moving images;
multiplication by $v_j(q)$ is not being asserted to preserve a fixed
image of the profile space.
There is no $\zeta$ dependence in $v_j$. Differentiating this complete
equation once more, including the time derivatives of its coefficients,
gives
\begin{equation}
 \partial_t^2\Psi_j=
 \frac{(\vartheta')^2}{16}\partial_\zeta^4\Psi_j
 -\frac{\vartheta'v_j}{2}\partial_\zeta^3\Psi_j
 +\left(v_j^2-\frac{\vartheta''}{4}\right)\partial_\zeta^2\Psi_j
 +a_j\partial_\zeta\Psi_j,
 \label{nfl:second}
\end{equation}
where the acceleration coefficient is, with all NS terms retained,
\begin{equation}
 a_j=\partial_tv_j=
 \left[\sum_i(\nu\Delta u_i-\partial_i p+f_i)\partial_i S_j
             +\sum_{i,l}u_i u_l\partial_i\partial_l S_j\right](t,X_t(q)).
 \label{nfl:accel}
\end{equation}
Indeed apply $\partial_t+u\cdot\nabla_x$ to
$\sum_i u_i\partial_iS_j$ and substitute \eqref{nfl:input}; the
derivative of $\partial_iS_j$ is exactly the displayed Hessian term.
In \eqref{nfl:second}, the cubic derivative occurs twice, each with
coefficient $-\vartheta'v_j/4$, proving its coefficient $-\vartheta'v_j/2$.

The full physical Laplacian has an exact material counterpart. Define
\[
 G_t=A^{-1}A^{-\mathsf T},\qquad
 \mathcal L_t=\sum_{a,b}\partial_{q_a}
                         (G_t^{ab}\partial_{q_b}).
\]
For every smooth scalar $b(x)$,
\begin{equation}
 \mathcal L_t(X_t^*b)=X_t^*(\Delta_xb).
 \label{nfl:laplace}
\end{equation}
Here is a proof retaining the first-order terms in $\mathcal L_t$.
The gradient chain rule gives $\nabla_qX_t^*b=A^{\mathsf T}
(\nabla_xb)\circ X_t$. For a compactly supported smooth test $\phi$,
integration by parts and the change of variables $x=X_t(q)$, whose
Jacobian is one, give
\[
 \int\phi\,\mathcal L_tX_t^*b\,dq
 =-\int\nabla_q\phi\cdot A^{-1}(\nabla_xb)\circ X_t\,dq
 =-\int\nabla_x(\phi\circ X_t^{-1})\cdot\nabla_xb\,dx
 =\int\phi\,X_t^*\Delta_xb\,dq.
\]
All functions are smooth, so the distributional equality is pointwise.
Apply this at fixed $\zeta$ to the full profile to obtain
\begin{equation}
 \mathcal L_t\Psi_j=(|\nabla S_j|^2\circ X_t)\partial_\zeta^2\Psi_j
                      +(\Delta S_j\circ X_t)\partial_\zeta\Psi_j.
 \label{nfl:profilelaplace}
\end{equation}
Consequently its complete advective viscous residual is
\[
 (\partial_t-\nu\mathcal L_t)\Psi_j
 =-\left(\frac{\vartheta'}4+\nu|\nabla S_j|^2\circ X_t\right)
                                  \partial_\zeta^2\Psi_j
   +(v_j-\nu\Delta S_j\circ X_t)\partial_\zeta\Psi_j.
\]
These formulas identify the two differential operators through an
explicit coordinate map and calculate their remaining coefficients;
they do not delete the viscosity or equate two time variables.

\subsection{Recovering the actual velocity from the arithmetic temporal defect}

For real $\theta$ the functions $Z_\theta,K_\theta$ and their spectral
derivatives are Schwartz on the real $s$ line, by the weighted Fourier
space proved above. With the original real kernel $k$ and multiplier
$m$, their Fourier representations are
\[
 Z_\theta(s)=4\int_\R k(\omega)e^{\theta\omega^2/4}e^{is\omega}\,d\omega,
 \qquad
 K_\theta(s)=4\int_\R m(\omega)k(\omega)e^{\theta\omega^2/4}
                                          e^{is\omega}\,d\omega.
\]
No source energy, derivative channel, or Fourier factor is changed.
Define $d_h(\theta)^2=\int_\R|h_\theta'(s)|^2\,ds$. Then
\begin{align}
 d_Z(\theta)^2&=32\pi\int_\R
                   \omega^2 k(\omega)^2e^{\theta\omega^2/2}\,d\omega,\notag\\
 d_K(\theta)^2&=32\pi\int_\R
        \omega^2|m(\omega)|^2 k(\omega)^2e^{\theta\omega^2/2}\,d\omega.
 \label{nfl:d}
\end{align}
For completeness, the factor follows without a Fourier convention
change. For a Schwartz amplitude $b$, insert $e^{-\epsilon s^2}$ in
the squared norm of $4\int b(\omega)e^{is\omega}d\omega$ and use Fubini.
The result is
\[
16\sqrt{\pi/\epsilon}\iint b(\omega)\overline{b(\eta)}
e^{-(\omega-\eta)^2/(4\epsilon)}d\omega d\eta.
\]
The Gaussian convolution, with mass $2\pi$, tends to $2\pi b$;
splitting at any fixed small distance and using continuity and the
Gaussian tail proves this convergence under the integral. Thus the
limit is $32\pi\int|b|^2$. On the original side dominated convergence
applies because the transform is Schwartz. Use respectively
$b=i\omega k e^{\theta\omega^2/4}$ and
$b=i\omega m k e^{\theta\omega^2/4}$ to prove \eqref{nfl:d}.

Both numbers are strictly positive. The previously proved $k(0)>0$
and continuity give an interval of nonzero $k$ containing nonzero
$\omega$. For the second integral the established full-real bound
$|m|\geq575/1752192$ applies. The weighted bounds also prove continuity
of these integrals in $\theta$ by domination on each compact parameter
interval. Hence for each compact interval $I\subset\R$ there are
finite constants $0<d_-\leq d_h(\theta)\leq d_+<\infty$ on $I$.

Define the actual arithmetic temporal defect by
\begin{equation}
 D_j(t,q,s)=\partial_t\Psi_j(t,q,s)
                   +\frac{\vartheta'}4\partial_s^2\Psi_j(t,q,s),
 \qquad s\in\R.
 \label{nfl:defect}
\end{equation}
Equation \eqref{nfl:first} gives $D_j=v_j h_{\vartheta}'(s+s_{j,t})$.
Real translation preserves Lebesgue measure, so
\begin{equation}
 \sum_j\|D_j(t,q,\cdot)\|_{L^2(\R)}^2
       =d_h(\vartheta(t))^2|v(t,q)|^2.
 \label{nfl:norm}
\end{equation}
Moreover the inverse is a literal integral followed by the original
matrix inverse:
\begin{equation}
 \begin{aligned}
 v_j&=\frac{\displaystyle\int_\R
      \overline{h_{\vartheta}'(s+s_{j,t})}D_j(t,q,s)\,ds}
                {d_h(\vartheta)^2},\\
 u(t,X_t(q))&=2W_1(X_t(q))v_1+W_2(X_t(q))v_2+W_3(X_t(q))v_3.
 \end{aligned}
 \label{nfl:recoveru}
\end{equation}
Substitution proves both inverse identities. The correspondence retains
$X_t(q)$ as part of both sides:
$(X_t,u\circ X_t)\leftrightarrow(X_t,D)$. The data $D$ alone are not
claimed to select a globally unique inverse sheet of the polynomial map.
At fixed $(t,q)$, the formula
identifies $\R^3$ with the three-dimensional real subspace of
$L^2(\R;\C)^3$ spanned componentwise by the three displayed shifted
profiles. It is not merely an analogy between their norms.

For a nonempty $K$, set
$M=\max_K\|J\|_{\mathrm{op}}$ and
$N=\max_K\|J^{-1}\|_{\mathrm{op}}$. They are finite and positive.
If $\vartheta([0,T))\subset I$, the exact map satisfies, for all $t<T$,
\begin{equation}
 \frac{d_-}{N}\|u(t)\|_{L^\infty(\R^3)}
 \leq \sup_q\left(\sum_j\|D_j(t,q,\cdot)\|_2^2\right)^{1/2}
 \leq d_+M\|u(t)\|_{L^\infty(\R^3)}.
 \label{nfl:blowupequiv}
\end{equation}
Indeed $|Ju|\leq M|u|$ and $|u|\leq N|Ju|$ on $K$, while both $u$
and $v$ vanish off $K$; \eqref{nfl:compact} preserves $K$ and
surjectivity preserves the supremum. Combine this with \eqref{nfl:norm}.
This proves equivalence of unbounded terminal velocity and unbounded
supremum of this specific arithmetic temporal-defect norm for every
input in the defined class. It supplies an exact detection map even
though no singular member is established here. At $u=0$ all these
defects vanish, as the same formulas require.
The full space--spectral energy has an exact identity too. Since $v=0$
off $K$, \eqref{nfl:norm} is integrable in $q$ and the volume-preserving
change of variables gives
\[
 \sum_j\int_{\R^3}\int_\R|D_j(t,q,s)|^2\,ds\,dq
   =d_h(\vartheta(t))^2\int_{\R^3}|J(x)u(t,x)|^2\,dx.
\]
The same constants therefore bound its square root between
$(d_-/N)\|u(t)\|_2$ and $d_+M\|u(t)\|_2$. In particular the map
transfers a uniform kinetic-energy bound into a uniform joint
$L^2(q,s)$ bound, while \eqref{nfl:blowupequiv} detects concentration
through the separate supremum in $q$. This conclusion follows from the
displayed integrals without a spatial truncation of the defect.
In particular $\vartheta(t)=0$ for every physical time is an explicit
allowed choice. Then the entire input stays equal to the actual
$\Xi$ (or $K_0$), and $D_j=\partial_t\Psi_j$; the constants in
\eqref{nfl:blowupequiv} use the single strictly positive value $d_h(0)$.

For an actual separated-link pair, the full coefficient is $a_{ef}K$.
The defect is then $a_{ef}D^K$, and its squared norm acquires exactly
$a_{ef}^2$; division in \eqref{nfl:recoveru} is available for
$a_{ef}\ne0$. If $a_{ef}=0$ this pair supplies the zero observation,
not a spurious inverse. For a pair with $a_{ef}>0$, recovery of $Z$ from
the profiles of $K$ additionally uses the already proved full-domain
inverse $\mathcal T_{\mathcal D}^{-1}$; its translation and derivative
commutation hold because multiplication by $m$ commutes with
$e^{is_{j,t}\omega}$ and $i\omega$ in the unchanged Fourier domain.

\subsection{Bounded profile values and the exact transported zero divisor}

For compact label set $L$, compact spectral set $E\subset\C$, and
bounded $\vartheta([0,T))\subset I$, \eqref{nfl:compact} puts every
argument $s_{j,t}(q)+\zeta$ in the compact set $S_j(K\cup L)+E$.
The jointly entire $Z$ and $K$ are bounded on that set times $I$.
Therefore
\[
 \sup_{t<T,\ q\in L,\ \zeta\in E}|\Psi_j(t,q,\zeta)|<\infty.
\]
The identical argument bounds every pure spectral derivative. It does
not bound time or spatial-label derivatives; \eqref{nfl:blowupequiv}
specifies a quantity that can grow even while these values stay bounded.

Let $h_\theta(s)$ have a zero at $z_*$ of multiplicity $m$ and write its
actual local factorization
$h_\theta(s)=(s-z_*)^m A(s)$, $A(z_*)\ne0$. The full inverse of
\eqref{nfl:skew} takes that zero to
\[
 \zeta_*(q)=z_*-s_{j,t}(q),\qquad
 \Psi_j=(\zeta-\zeta_*(q))^m A(s_{j,t}(q)+\zeta).
\]
Thus it preserves multiplicity and the whole analytic unit on every
complex fibre. Its complete jet formula is
\[
 \partial_\zeta^{m+n}\Psi_j(t,q,\zeta_*(q))
                 =\frac{(m+n)!}{n!}A^{(n)}(z_*),\qquad n\geq0,
\]
with all lower derivatives zero. This follows by multiplying the
Taylor series of the unchanged $A$ by $(\zeta-\zeta_*)^m$; no unit
coefficient has been discarded. Since $s_{j,t}(q)$ is real,
$\operatorname{Im}\zeta_*(q)=\operatorname{Im}z_*$. At Newman time
$\vartheta=0$ with $h=Z$, the full map to the completed-zeta argument is
\[
 \rho=\tfrac12+i(s_{j,t}(q)+\zeta),\qquad
 \zeta=-i(\rho-\tfrac12)-s_{j,t}(q),\qquad
 \operatorname{Re}\rho=\tfrac12-\operatorname{Im}\zeta.
\]
Consequently this NS-induced real translation preserves each zero's
distance from the critical line exactly; the full arithmetic time
evolution is still the specified $Z_\vartheta$. For $h=K$ the same
coordinate proof transports the zeros of the four-term profile; these
are converted to a statement about $Z$ only by actually applying the
proved inverse operator to the whole profile. The inverse does not
identify individual zeros of two different entire functions. These
explicit maps give both the velocity detection and the precise zero
transport of the construction, without deducing an off-critical zero
from growth of the temporal defect.
